\section{The certified decision}\label{sec:dec}

An \emph{enclosure} of $\log r(x,y)$ is an interval $[\ell_-,\ell_+]$ that contains it; Algorithm~\ref{alg:row}
builds it from the four density enclosures of Theorem~\ref{thm:dens} (and Proposition~\ref{prop:left}) and the
rounding of their sum (Appendix~\ref{app:round}). Enclosures are intersected across levels, so they never widen.

\begin{theorem}[Decisions taken are exact]\label{thm:dec}
If Algorithm~\ref{alg:row} returns, it returns what the exact kernel returns for the same $y$ and the same
uniform $U$ (the limit of its bits). At a level whose enclosure has half-width $E$, the probability, given $x$
and $y$, that the decision is not taken is at most $2E+2^{1-\beta}$, with $\beta$ the number of bits of $U$
drawn.
\end{theorem}
\begin{proof}
The exact kernel moves iff $U<\min(1,r)$. Accept branch: if $\ell_-\ge0$ then $r\ge1>U$; if $\ell_-\ge\log U_+$
then $U<U_+\le e^{\ell_-}\le r$. Reject branch: $\ell_+\le\log U_-<0$ gives $r\le e^{\ell_+}\le U_-\le U$ and
$r<1$. The undecided event is contained in $\{\log U_+>\ell_-\}\cap\{\log U_-<\min(0,\ell_+)\}$; for exact $U$ its
probability is at most $e^{\min(0,\ell_+)}-e^{\min(0,\ell_-)}\le2E$ as $e^t$ is 1-Lipschitz on $t\le0$, and the
bit interval adds at most its width on each side.
\end{proof}

\subsection{The precision ladder}\label{sec:ladder}

\begin{proposition}[Termination with a precision ladder]\label{prop:ladder}
Let level $j=0,1,\dots$ evaluate the four log-densities at working precision $p_j$ bits ($p_0=53$, $p_1=106$,
$p_j\to\infty$), with $h_j$ and the truncation chosen so that $E_{\mathrm d}+E_{\mathrm t}\le2^{-p_j}$, and hold
$\beta_j\ge p_j$ bits of $U$. Suppose every function the evaluator uses ($+,-,\times,/$, real and complex square
root, $\exp$, $\mathrm{expm1}$, $\log$, $\mathrm{log1p}$, $\sin$, $\cos$, and for the majorants $\sinh$) is
available at precision $p_j$ with a rigorous error bound proportional to $2^{-p_j}$ (correctly rounded MPFR and
MPC, or Arb). Then the half-widths satisfy $E_j\le C(x,y)\,p_j^{3/2}\,2^{-p_j}$, the number of levels $J$ satisfies
$P(J>j\mid x,y)\le2E_j+2^{1-\beta_j}\to0$, so $J<\infty$ with probability one and the implemented kernel equals
the exact kernel $P$ (under (A0) for the draws). If level $j$ costs at most a polynomial in $p_j$, the expected
cost $\sum_j\mathrm{cost}_jP(J\ge j)$ is finite and dominated by level 0.
\end{proposition}
\begin{proof}
At fixed $(x,y)$ the discretization bound decays like $e^{-2\pi a/h}$, so $h_j$ is of order $1/p_j$; the
integrand decays like $e^{-2xy^2}$, so the truncation point is of order $\sqrt{p_j}$; the node count $N_j$ is
therefore of order $p_j^{3/2}$. Ball arithmetic bounds the rounding of a recursive sum of $N_j$ terms by a constant
times $N_j2^{-p_j}$ times a quantity that converges as $p_j$ grows (the sum of the moduli of the terms); add the
mathematical part $2^{-p_j}$. Any polynomial in $p_j$ in place of $p_j^{3/2}$ gives the rest.
Theorem~\ref{thm:dec} at each level; $P(J=\infty)\le\inf_j(2E_j+2^{1-\beta_j})=0$; on $J<\infty$ the outputs agree
by Theorem~\ref{thm:dec}. Cost: $P(J\ge j)$ falls geometrically in $p_{j-1}$ up to a polynomial factor, while
$\mathrm{cost}_j$ grows polynomially.
\end{proof}

The ladder is unbounded: $h_j\to0$, $N_j\to\infty$ and $p_j\to\infty$. A build that caps any of them, as
\texttt{arbpg.py} caps $h$ at $0.75\cdot2^{-39}$ and the node count at 20000, keeps a positive floor even under
exact arithmetic wherever a cap binds before $E_{\mathrm d}+E_{\mathrm t}$ is small; it is then a build that stops at a
floor (Proposition~\ref{prop:floor}), with a floor no pin reached. Remark~\ref{rem:coverage}'s ``refining $h$
alone drives the enclosure to zero'' is likewise a statement about the uncapped ladder.

What it takes, in numbers (Section~\ref{sec:floor}, at the pins, with the lead in double-double): at level 0 (53 bits, tilted form) each log-density's
enclosure has radius at most \FLOORDENS{} at every pin, and level 1 is reached with probability at most about
\FLOORSTEP{} per row step; at 106 bits the radius is at most \FLOORHI{}
per density, so level 2 is reached with probability of order \FLOORHISTEP. Level 1 needs double-double versions of the
node formula of Proposition~\ref{prop:cf} (the complex functions above at 106 bits); the lead is in double-double
at every level already; levels 2 and up need an
arbitrary-precision library with rigorous functions, a new system dependency (MPFR and MPC) or a bundled one; the
majorants $E_{\mathrm d},E_{\mathrm t}$ may stay in double, since they are driven below $2^{-p_j}$ by $h$ and $N$,
not by precision, as long as their own rounding is bounded rigorously. The contour and the saddle stay in double
at every level (any choice is valid).

\begin{proposition}[A build that stops at a floor]\label{prop:floor}
Suppose instead the build stops at level $j_f$ with half-width at most $E_f$, takes a fixed decision there and
tallies the event. Then $\TV(\widetilde P(x,\cdot),P(x,\cdot))\le2E_f+2^{1-\beta}$ for every $x$; a sweep of $n$
rows differs from the exact sweep by at most $n(2E_f+2^{1-\beta})$; and $N$ sweeps from the same start by at most
$Nn(2E_f+2^{1-\beta})$, with no ergodicity assumption.
\end{proposition}
\begin{proof}
The two kernels coincide off the floor event (Theorem~\ref{thm:dec}). Run both chains on the same randomness; they
agree until the first floor event, whose probability over $N$ sweeps is at most the sum of the per-row bounds
(also \cite[eqs.~(3.3)--(3.4)]{mitro}, which use only $\tau(P^i)\le1$).
\end{proof}

\begin{proposition}[A fixed top level is not exact]\label{prop:fixedtop}
If the top level has positive half-width on a set of $(x,y)$ of positive $\pi\otimes q$ measure (as it has at any
finite precision), the kernel with a fixed decision at the floor differs from $P$, and $\pi$ is not invariant for
it in general. A tally counts floor events; it does not remove them.
\end{proposition}
\begin{proof}
With the fixed decision ``keep $x$'', $\widetilde P(x,\{x\})-P(x,\{x\})=\int q(y)P(e^{\ell_-}\le U<\min(1,r)\mid x,y)\,dy$,
positive wherever the enclosure has positive width with $r$ inside it; so $\widetilde P\ne P$. That $\pi$ is then
not invariant in general is shown on two states, the argument of the second referee. Let the target be
$(\pi_0,\pi_1)$ and the proposal $(q_0,q_1)$, with $r=r(0,1)=\pi_1q_0/(\pi_0q_1)<1$, and let the half-width $E$ at
the top level be positive and below $\log(1/r)$. From 1 the ratio is $1/r>1$, its enclosure lies above 0, and the
move to 0 is always decided: the flow $\pi_1q_0$ from 1 to 0 is the exact kernel's. From 0 the move is taken
only when $U$ is certainly below $r$, which removes the probability $\epsilon>0$ that $U$ falls in the band below
$r$: the flow from 0 to 1 is $\pi_0q_1(r-\epsilon)=\pi_1q_0-\pi_0q_1\epsilon$. On two states $\pi$ is invariant iff
the two flows are equal, so it is not. Toy~6 (Appendix~\ref{app:toys}) gives the exact numbers on two and three
states. The claim is non-invariance in general, which these settle; by how much a given row's target moves at a
floor is bounded by Theorem~\ref{thm:mitro}, not computed.
\end{proof}

\begin{theorem}[Stationary error at a floor]\label{thm:mitro}
At fixed $(b,\psi)$, if the floor chain $\widetilde P$ of one row has an invariant distribution $\tilde\pi$, then
$\TV(\tilde\pi,\PG(b,\psi))\le(\lceil M^*\log2\rceil+M^*)\,\bar E$, where $\bar E=2E_f+2^{1-\beta}$ bounds the
floor probability from every $x$.
\end{theorem}
\begin{proof}
\cite[Cor.~3.1, eq.~(3.8)]{mitro}: if $X$ is uniformly ergodic, $\|e_xP^n-\pi\|\le C\rho^n$ for all $x$ (their
(2.2)), and $\widetilde X$ has an invariant measure $\tilde\pi$, then $\|\tilde\pi-\pi\|\le(\hat n+C\rho^{\hat
n}/(1-\rho))\|E\|$ with $E=\widetilde P-P$ and $\hat n=\lceil\log_\rho C^{-1}\rceil$. Their norm is the total
variation of a signed measure, twice $\TV$, and $\|E\|=\sup_x\|e_xE\|$. Theorem~\ref{thm:row} gives (2.2) with
$C=2$, $\rho=1-1/M^*$; Proposition~\ref{prop:floor} gives $\|E\|\le2\bar E$. Then $C\rho^{\hat n}\le1$,
$1/(1-\rho)=M^*$, and $\hat n\le\lceil M^*\log2\rceil$ since $-\log(1-1/M^*)\ge1/M^*$; halve to pass from the
norm to $\TV$.
\end{proof}
For the whole sampler no such stationary bound is available (uniform ergodicity of the BART chain is not known);
Proposition~\ref{prop:floor}'s finite-horizon bound is what holds, and the tally measures it.

\begin{proposition}[A fresh draw on an undecided decision]\label{prop:fresh}
Replacing the step's output by a fresh exact draw from $\PG(b,\psi)$ whenever the decision is undecided does not
keep $\PG(b,\psi)$ in general.
\end{proposition}
\begin{proof}
Counterexample, in exact rational arithmetic (Toy 2, Appendix~\ref{app:toys}): three states, target
$(\frac12,\frac13,\frac16)$, uniform proposal, undecided band $|U-a(x,y)|<(1+x+2y)/100$ resolved by a fresh draw:
$\max_j|(\pi\widetilde P)_j-\pi_j|=1/120$ and the stationary distribution is $0.0108$ from $\pi$ in $\TV$.
The mechanism, to first order in a relative band ($|\log U-\log r|<E$): the undecided probability from $x$ to $y$
is $2E\,r(x,y)$ when $w(y)<w(x)$ and $O(E^2)$ in total otherwise, the exact kernel then moves with conditional
probability $\frac12+O(E)$, and the change of $\pi\widetilde P-\pi$ has density
$E\,f(z)\,[2c-Q(w>w(z))-F(w<w(z))/w(z)]+O(E^2)$, $c$ a constant, $Q,F$ the proposal and target measures. The
bracket tends to $2c-1$ where $w(z)\to\inf w=0$ and to $2c-1/M^*$ where $w(z)\to M^*$, so it is constant only if
$M^*=1$, i.e. $f=q$. Whether a decision resolves depends on the current value, which is why the fresh draw
biases.
\end{proof}

\section{Floating point: the standard every sampler in the package meets}\label{sec:ideal}

\paragraph{The idealization.} Every exactness claim about a Monte Carlo sampler, here and in the literature, is
about the idealized algorithm: uniforms are exact real uniforms and arithmetic is exact (the model of
\cite[Ch.~I.1]{devroye}). No floating-point sampler meets it in total variation: its output takes finitely many
values, so its distribution is at total variation 1 from any continuous one. The meaningful comparison is a
coupling: drive the idealized algorithm and the engine with the same random stream, the engine's uniform
$\hat U$ being the idealized $U$ on the engine's grid. Rounding then moves continuous outputs by a few units in the
last place, which no test can see, and occasionally flips a discrete decision (accept or reject, which branch),
which changes the output. The probability of a flip per draw is the floor.

\begin{lemma}\label{lem:flip}
A comparison of a uniform against a threshold $c\in[0,1]$, computed as $\hat c$ with $|\hat c-c|\le\epsilon$,
using a uniform on a grid of spacing $g$, decides differently from the idealized comparison with probability at
most $g+\epsilon$.
\end{lemma}
\begin{proof}
The decisions differ only if $U$ lies between $c$ and $\hat c$, or within $g$ of the threshold on the grid.
\end{proof}

\paragraph{Today's exact sum.} The engine's uniform is 32-bit ($g=2^{-32}=2.3\times10^{-10}$; Mersenne Twister,
\texttt{randomBase.c}). Its $\PG(1,z)$ sampler (\texttt{ext\_rng\_simulatePolyaGamma}, the alternating-series
sampler of \cite[Sec.~4.1]{psw}) compares, per proposal, a fresh uniform with the probability of the exponential
piece (computed through \texttt{pnorm} in double), and a fresh uniform times the first series coefficient with
partial sums $S_n$ computed in double; the truncated inverse-Gaussian piece adds one or more uniform comparisons
of its own. By Lemma~\ref{lem:flip} each comparison can flip with probability up to about $2^{-32}$, so a
$\PG(1)$ draw has a floor of about $5\times10^{-10}$ to $10^{-9}$, and $\PG(b)$, the sum of $b$ such draws, of
order $10^{-9}b$ per row per sweep: a few times $10^{-8}$ at $b=64$, about $10^{-6}$ at $b=10^3$. Its gamma draws
(Ahrens-Dieter GD, as in R) evaluate their acceptance tests with polynomial approximations whose coefficients
carry seven digits, a further floor of the same kind, not quantified here. None of this is tallied, and the sum
is exact in the only sense any sampler is: under (A0).

\paragraph{The step under the same coupling.} A row step makes $K+1$ gamma draws (GD makes one or more
uniform comparisons per draw, about two counting its exponential and normal parts), one acceptance comparison
against a uniform of at least 53 bits, and at the floor a fixed decision, so its floor is about
$2(K+1)\,2^{-32}+\bar E$ with $\bar E=2E_f+2^{1-\beta}$ the floor probability of the certified decision. In the
tilted form $\bar E$ is at most \FLOORSTEP{} at every pin
(Section~\ref{sec:floor}; at the pins, in ball arithmetic, with the lead in double-double), below the
$2^{-32}=2.3\times10^{-10}$ of a single comparison; in the untilted form it
reached $10^{-7}$ at $b=10^5$, $|\psi|=150$. The gamma part, $2(K+1)2^{-32}$, is below the sum's $2b\,2^{-32}$ or
more wherever $K<b$, which holds at every whole-number shape the rule sends to the step except at large $|\psi|$
with $b$ below a few hundred (at $|\psi|=150$, $K=140$ at $b=100$); there the two are of the same order.

\paragraph{Same standard, or stricter?} Two things differ. (1) The step's error enters a chain: a per-step
error $\bar E$ moves the row's stationary distribution by up to $(\lceil M^*\log2\rceil+M^*)\bar E$
(Theorem~\ref{thm:mitro}), where a fresh draw's moves it by its own floor; the factor is of the size of $M^*$
(between 1.0 and 36 where it is certified from below; no upper bound is certified). (2) Whether a floor event
happens depends on the current value, as in Proposition~\ref{prop:fresh}; but so does every rounding flip in every
sampler, whose thresholds depend on the state. Neither is a different kind of error: (1) is a constant factor and
(2) is shared. So the step can be held to the package's standard: exact under (A0), with a floor that is quoted,
here with the factor of (1). Under exact arithmetic the rounding radius is zero and refining $h$ alone drives the
enclosure to zero (Remark~\ref{rem:coverage}), so the decision always resolves; the floor is the rounding part
only, as for every other sampler. With the tilted form, the pins' \FLOORSTEP{} times the factor stays below one
uniform's grid as long as $M^*$ is below about 20 (and the floor at most shapes is far smaller); no upper bound on
$M^*$ is certified, and the engine's floor is still to be measured. With the untilted form the floor alone was up to $10^{-7}$. A stricter standard (exact in the implemented arithmetic) cannot be met by the step
alone, since its proposal is drawn by floating-point gamma generators; Section~\ref{sec:maint} puts the choice to
the maintainer.

\section{The sweep}\label{sec:sweep}

Let $\theta$ denote every unknown other than the latents (trees and leaves, the combiner's glue, $k$, DART's
weights, nbinom's $r$), and $\pi(\theta,\omega)$ the posterior given the data the larger sampler installed.

\begin{lemma}[Conditional independence]\label{lem:condind}
$\pi(\omega\mid\theta)=\prod_{i\ \mathrm{active}}\PG(\omega_i\mid b_i(\theta),\psi_i(\theta))$.
\end{lemma}
\begin{proof}
\cite[Thm.~1]{psw}: $(e^{\psi})^{y}/(1+e^{\psi})^{b}=2^{-b}e^{\kappa\psi}\int_0^\infty e^{-\omega\psi^2/2}
p(\omega\mid b,0)\,d\omega$, $\kappa=y-b/2$, for $b>0$. The augmented posterior is the prior times $\prod_i
2^{-b_i}e^{\kappa_i\psi_i}e^{-\omega_i\psi_i^2/2}p(\omega_i\mid b_i,0)$, which given $\theta$ factorizes into the
tilts $\PG(b_i,\psi_i)$. For nbinom the row's likelihood is $\mu_i^{y_i}r^r/(\mu_i+r)^{y_i+r}$ up to factors
free of $\eta_i$, that is $(e^{\psi_i})^{y_i}/(1+e^{\psi_i})^{y_i+r}$ with $\psi_i=\eta_i-\log r$ ($\eta_i=\log
\mu_i$, offset included) and $b_i=y_i+r$, as the engine forms it (\texttt{NBResponse::anchor}); for binomial $b_i$
is the trials times the row weight, with the weight entering as a power of the row's likelihood.
\end{proof}

\begin{theorem}[Sweep invariance]\label{thm:sweep}
Let the latent step act as $L(\theta,\omega;d\omega')=\prod_iR_i(\omega_i,d\omega_i')$, where each $R_i$ leaves
$\PG(b_i,\psi_i)$ invariant (a fresh draw or the step), chosen by anything except $\omega$ (marks, shape,
settings). Let every other step of the sweep leave $\pi$ invariant. Then the sweep leaves $\pi$ invariant.
\end{theorem}
\begin{proof}
For a set $B$ with sections $B_\theta$,
\[\int\pi(d\theta,d\omega)L(\theta,\omega;B_\theta)=\int\pi(d\theta)\int\prod_i\PG_i(d\omega_i)\,
\Bigl(\prod_iR_i\Bigr)(\omega,B_\theta)=\int\pi(d\theta)\Bigl(\prod_i\PG_i\Bigr)(B_\theta),\]
by
Lemma~\ref{lem:condind} and invariance of each factor on the product measure; this is $\pi(B)$. A composition
of invariant kernels is invariant. Rows taken one at a time give the same result, each $R_i$ being invariant for
$\pi(\omega_i\mid\theta,\omega_{-i})=\PG_i$.
\end{proof}

\begin{proposition}[Which later steps read $\omega$]\label{prop:steps}
After the latent step, the trees read $\omega$ (through the working response $\kappa_i/\omega_i$ and weights
$\omega_i$) and move $\psi$; each is a move for $\pi(\text{trees}\mid\omega,\text{rest})$, so after it
$(\theta,\omega)$ is again distributed as $\pi$, and the next sweep's step meets $\omega$ from its conditional at
the new $\psi$. The combiner's glue (where present) also reads $\omega$ through the working response. The level
shift moves constants between trees with every row's total fit unchanged, so no $\psi_i$ changes and it reads no
latent. The probit rescaling runs only for the probit family. $k$ and DART read the leaves and split counts and
change no $\psi_i$; $\omega$ is conditionally independent of them given the leaves, so their draws are full
conditionals. The only step that changes a row's conditional without reading $\omega$ is the nbinom $r$ step
(Section~\ref{sec:nb}), plus, between sweeps, the setters (Section~\ref{sec:setters}).
\end{proposition}
\begin{proof}
Each statement is the role of the step in the sweep order of the latents-first plan; invariance of each step for
$\pi$ is the standing property of today's sampler, unchanged by this design. The tree step is a Metropolis or Gibbs
move for its conditional given $\omega$, hence invariant for the joint.
\end{proof}

\section{Negative binomial with an estimated dispersion}\label{sec:nb}

\begin{theorem}\label{thm:nb}
Let $K_\theta(r,dr')$ leave $\pi(r\mid\theta_{-r})$ ($\omega$ integrated out) invariant and not read $\omega$; the
build's grid draw is $K_\theta(r,\cdot)=\pi(\cdot\mid\theta_{-r})$. Follow it by: if $r'\ne r$, a fresh exact draw
$\omega_i\sim\PG(y_i+r',\psi_i)$ at every active row; if $r'=r$, any latent step as in Theorem~\ref{thm:sweep}.
Here $\omega_i$'s conditional is $\PG(y_i+r',\eta_i-\log r')$: both its shape and its tilt move with $r$. This
block leaves $\pi$ invariant, and the per-row choice, made from $b_i=y_i+r'$, does not depend on $\omega$.
\end{theorem}
\begin{proof}
Fix $\theta_{-r}$ and start from $(r,\omega)\sim\pi(r\mid\theta_{-r})\pi(\omega\mid r,\theta_{-r})$. Since
$K$ does not read $\omega$, $\omega\perp r'\mid(r,\theta_{-r})$, so given $(r,r')$ with $r'=r=\rho_0$, $\omega\sim
\pi(\omega\mid\rho_0,\theta_{-r})$, and the step keeps that. On $r'\ne r$, the fresh draw has that distribution
at $\rho_0=r'$ regardless of the input. So given $(r,r')$, $\omega'\sim\pi(\cdot\mid r',\theta_{-r})$, and
$P(r'\in C)=\int\pi(dr\mid\theta_{-r})K(r,C)=\pi(C\mid\theta_{-r})$. Hence $(r',\omega')\sim\pi(r,\omega\mid
\theta_{-r})$, and the block, holding $\theta_{-r}$ fixed, is invariant for $\pi$. Positive probability of
$\{r'=r\}$ is not needed. All of $b_i$, $\kappa_i$ and $\psi_i=\eta_i-\log r$ involve $r$; on
$\{r'=r\}$ none changes, and on $\{r'\ne r\}$ the fresh draw is at $(y_i+r',\eta_i-\log r')$. The grid has whole-number values and nbinom takes no row weights, so every $b_i$ is a whole number and the
fresh draw exists.
\end{proof}

Applying the step after $r$ moved is not invariant: in exact arithmetic (Toy 1) the stationary distribution is
$0.039$ from $\pi$ in $\TV$, while the designed rule is exact. This is the situation of \cite[Sec.~3, Sampler
Fragment 1 and the remark after it]{vdj}: ``an MH update in a step that follows a reduced step is
problematic''; the fresh draw is their remedy, a direct draw after the reduced step. A real-valued dispersion
makes $b_i$ fractional when $r$ moves, with no fresh draw; an $r$ step that reads $\omega$ (invariant for
$\pi(r\mid\omega,\theta_{-r})$) needs none, and then the step is valid at every row.

\section{The per-row rule}\label{sec:rule}

\begin{theorem}\label{thm:rule}
Let row $i$ use the fresh draw when $\chi_i=1$ and the step when $\chi_i=0$, with $\chi_i$ any measurable function
of $(\theta,\omega_{-i})$, the marks, the settings and randomness independent of $\omega_i$; dec-B590's rule is
$\chi_i=1$ iff $b_i\le64$ and $b_i\in\mathbb Z$, or the row is marked. Then the row's update leaves
$\PG(b_i,\psi_i)$ invariant. A $\chi_i$ read from $\omega_i$ can break invariance.
\end{theorem}
\begin{proof}
Given the conditioning variables, $\chi_i$ is fixed and both kernels are invariant. Counterexample for a choice
read from $\omega_i$ (Toy 5): on four states, the step at states 0, 1 and a fresh draw at 2, 3 gives
$\max|\pi\widetilde P-\pi|=0.167$ and stationary $\TV$ $0.258$, though each kernel alone is exact.
\end{proof}

\paragraph{The cancel.} Whether a row is processed before a cancel lands is a choice of the kind the theorem
excludes when the time a row takes depends on its own $\omega_i$ (the refinement does), so a cancel polled inside a
row's step leaves that row unchanged by a rule read from $\omega_i$. Today's refresh polls once every $2^{16}$
$\PG(1)$ draws (\texttt{RefreshPoll}), a count that can fall inside a row's sum, and a stopped row is left
unchanged; the count does not depend on the values drawn, only the moment the cancel arrives does. For the step,
the build should poll between rows only (counting rows or steps), never inside a row's refinement, whose length
depends on $\omega_i$. Even
then a cancel ends the run partway through a sweep, and the state it leaves is a start for a later run
(Section~\ref{sec:starts}), not a draw from $\pi$; that is today's behavior too.

\section{Setters, marked rows and the two contracts}\label{sec:setters}

The code that runs the BART sampler inside a larger sampler has its own unknowns $\phi$ and installs data from them
(offsets, weights, responses, predictors) through setters between sweeps; a setter marks every row whose
$(b_i,\psi_i)$ or activity it may change, and setData and setState mark every row. The joint target is
$\Pi(\phi,\theta,\omega)=p(\phi)\,\pi_\phi(\theta,\omega)$, where $p(\phi)$ is $\phi$'s marginal under $\Pi$ (its
prior times the marginal likelihood of the BART part, normalized), not its prior.

\begin{definition}[Contracts]
Under \emph{contract C} the larger sampler's step on $\phi$ leaves $\Pi(\phi\mid\theta,\omega)$ invariant; it may
read $\omega$ (through getLatents). Under \emph{contract I} its step leaves $\Pi(\phi\mid\theta)$ invariant and
does not read $\omega$ (the latents are integrated out).
\end{definition}

\begin{theorem}\label{thm:setters}
(a) Under contract C the composite chain leaves $\Pi$ invariant, with fresh draws or the step at marked rows.
(b) Under contract I it leaves $\Pi$ invariant when every marked row is drawn fresh at the next sweep's start and
the marks cover every row whose conditional changed. (c) Under contract I the step at a marked row whose conditional changed does not keep
$\Pi$ in general. This includes every row a setter or setState marks when its shape is fractional: the design
serves them by the step from the installed value, and under contract I that is the case of (c) whenever the call,
made between sweeps of a chain, changes such a row's $(b_i,\psi_i)$ or its activity. A call that changes no row's
conditional (the same data handed again, or a state saved under the data now installed) leaves the step exact
(Section~\ref{sec:maint}, Question~2, gives the options and the ruling).
\end{theorem}
\begin{proof}
(a) The $\phi$ step is invariant for $\Pi$ (a move for a full conditional); the sweep is invariant for
$\pi_\phi$ at each $\phi$ (Theorem~\ref{thm:sweep}) whatever the marks choose. (b) From $(\phi,\theta,\omega)\sim
\Pi$, $\phi'\sim K_\theta(\phi,\cdot)$ not reading $\omega$, so $\omega\mid(\phi,\phi',\theta)\sim\prod_i\PG(b_i(
\phi),\psi_i(\phi,\theta))$. Unmarked rows have the same factor at $\phi'$; marked rows are replaced by draws from
their factor at $\phi'$; by conditional independence, $\omega'\mid(\phi,\phi',\theta)\sim\pi_{\phi'}(\omega\mid
\theta)$, while $(\phi',\theta)$ has the $\Pi$ marginal by invariance of $K$. So $(\phi',\theta,\omega')\sim\Pi$;
then the rest of the sweep. (c) Toy 4 in exact arithmetic: under contract I the step at the changed row leaves a
stationary distribution $0.0154$ from $\Pi$ in $\TV$ ($\max|\Pi\widetilde P-\Pi|=0.00323$); the fresh draw and
both variants under contract C are exact.
\end{proof}

\paragraph{Scope: a fixed set of rows.} Part (a) holds for a fixed set of rows; an activity change at a
fractional shape is covered only on the pseudo-prior extension. When the active set is among the larger sampler's
unknowns, $\Pi$ defines $\omega$ only at the active rows, so a move that brings row $j$ in cannot condition on
$\omega_j$; an inactive row keeps its last $\omega_j$, stale while the trees move. On the pseudo-prior extension
(B.~P. Carlin and S.~Chib, 1995) each inactive row carries $\omega_j\sim\PG(b_j,0)$, $b_j$ its shape when
active, which does not depend on the trees: the tree step and the frozen $\omega_j$ leave the extension invariant,
its marginal is $\Pi$, and an activity update from its conditional given every row's $\omega$, the stale ones
included, is a contract-C step on it. Two independent referees proved this; their reports are beside this note. A
row that setData adds, or a weight moved from 0 to a value the larger sampler draws, starts from a value the step
forgets only over sweeps, and is exact under neither contract until then. At whole-number shapes a returning row
is drawn fresh and is exact either way.

\paragraph{Outer steps under neither contract.} The two contracts do not cover every step a larger sampler can
take. A step that reads $\omega$ (in its proposal, say) but is invariant only for $\Pi(\phi\mid\theta)$, or that
mixes steps of both kinds, is covered by neither: after it, $\omega$ and $\phi'$ are dependent given $\theta$ in a
way neither $\Pi(\phi\mid\theta,\omega)$ nor the latent conditional accounts for, and fresh draws at marked rows do
not repair the unmarked rows, whose $\omega_i$ is now correlated with $\phi'$. Nothing is guaranteed for such a
step, with the step or with today's fresh draws: Theorem~\ref{thm:setters} gives exactness for an outer step of
kind C or of kind I only. The engine's documentation of getLatents and the setters should say so: reading the
latents commits the caller to contract C (a move for $\phi$'s full conditional given $\omega$); not reading them,
with a move for $\phi$ given $\theta$ alone, is contract I.

\begin{lemma}[Only a fresh draw forgets the old conditional]\label{lem:complete}
Let $K$ be a Markov kernel on $(0,\infty)$ with $\int\PG(b,\psi_0)(dx)K(x,A)=\PG(b,\psi_1)(A)$ for every Borel
$A$ and every $\psi_0$ with $\psi_0^2/2$ in an open interval $I$. Then $K(x,\cdot)=\PG(b,\psi_1)$ for almost every
$x$.
\end{lemma}
\begin{proof}
For fixed $A$, $g(x)=K(x,A)-\PG(b,\psi_1)(A)$ is bounded and $\int g(x)\fz(x)e^{-\sigma x}dx=0$ for $\sigma\in I$.
The left side is the Laplace transform of the finite signed measure $g\fz\,dx$, analytic for $\mathrm{Re}\,\sigma>
-c_1$; vanishing on $I$, it vanishes there identically, so $g\fz=0$ a.e. by uniqueness, and $\fz>0$. Apply this to
the intervals with rational endpoints, a countable class closed under intersection that generates the Borel sets.
\end{proof}

\begin{corollary}[Fractional-shape marked rows]\label{cor:frac}
Let a marked row's kernel $K_{\psi_1}$ depend on the new $\psi_1$ and the row's data but not on the old $\psi_0$
(the step keeps no record of it). If the chain leaves $\Pi$ invariant for every contract-I outer step, then
$K_{\psi_1}(x,\cdot)=\PG(b,\psi_1)$ for almost every $x$: it is an exact fresh draw. At fractional $b$ none is available here: none is published below 1, and the sampler of
\cite[Sec.~4]{wps} for non-integer shapes rests on their Conjecture 8, of which they write ``We numerically verify this is the case for $h\in[1,4]$''.
Kernels that use $\psi_0$ exist in principle, such as the transport $x\mapsto F_{\psi_1}^{-1}(F_{\psi_0}(x))$ of
the distribution functions, or rejection from $q$ given a certified $M^*$; neither is built. So, at fractional
shapes, no kernel built here keeps every contract-I caller exact after a mid-chain change; contract-C callers are
exact with the step. The design takes every setter at fractional shapes (Section~\ref{sec:maint}, Question~2),
and its help tells a contract-I caller to keep whole-number weights on the rows it changes.
\end{corollary}
\begin{proof}
Under contract I, $\omega_i\mid(\phi,\phi',\theta)\sim\PG(b,\psi_0)$ with $\psi_0=\psi_i(\phi,\theta)$, so
invariance needs $\int\mu(d\psi_0)\int\PG(b,\psi_0)(dx)K_{\psi_1}(x,\cdot)=\PG(b,\psi_1)$, where $\mu$ is the
conditional of $\psi_0$ given $(\phi',\theta)$: one condition on a mixture, which by itself does not force
$K_{\psi_1}$ to be a fresh draw. Necessity comes from asking it of every contract-I outer step, including
deterministic ones. Let $\phi$ be real (an offset shift, say) with a continuous conditional distribution function
$F_\theta$ under $\Pi(\phi\mid\theta)$, strictly increasing on an interval where $\psi_i(\cdot,\theta)$ is continuous
and not constant, and for $s\in(0,1)$ let $T_s(\phi)=F_\theta^{-1}(F_\theta(\phi)+s\bmod1)$, the quantile rotation
(the referees' construction). If $\phi\sim\Pi(\cdot\mid\theta)$ then $F_\theta(\phi)$ is uniform on $(0,1)$, which
the rotation by $s$ preserves, so $T_s\phi\sim\Pi(\cdot\mid\theta)$: each $T_s$ is a deterministic contract-I step,
in every model with a continuous conditional. Given $\phi'=T_s\phi$ it fixes $\phi=T_s^{-1}\phi'$ almost surely, so
$\mu$ is a point mass at $\psi_0=\psi_i(T_s^{-1}\phi',\theta)$, while $\psi_1=\psi_i(\phi',\theta)$ does not depend
on $s$. A kernel that ignores which step was taken must meet the condition for every $s$; for $s$ in a countable
dense set it holds for almost every $\phi'$, and as $s$ varies $\psi_0$ sweeps an interval, so $\psi_0^2/2$ sweeps an
open interval $I$ on a dense subset of which the Laplace transform of Lemma~\ref{lem:complete}'s proof vanishes,
hence on all of $I$ by continuity. That is the hypothesis of Lemma~\ref{lem:complete}, which gives the claim. For a single outer step, a kernel tuned to its $\mu$ is not
excluded, but it would have to know $\mu$, which the engine does not.
\end{proof}

\section{Starting values}\label{sec:starts}

Let $S_F$ be the sweep with fresh draws at every row (today's sampler) and $S_R$ the sweep with the step at the rows
an $\omega$-free rule picks.

\begin{theorem}\label{thm:starts}
(a) $S_R(\xi,\cdot)\ge\epsilon(\theta)S_F(\xi,\cdot)$ with $\epsilon(\theta)=\prod_{i\ \mathrm{step}}1/M^*(b_i,
\psi_i)>0$. (b) If $S_F$ is $\varphi$-irreducible, so is $S_R$; if $S_F$ is also aperiodic, so is $S_R$. (c) Then,
by \cite[Thm.~4]{rr}, $\|S_R^n(\xi,\cdot)-\pi\|\to0$ for $\pi$-almost every start $\xi$, and for every initial
distribution $\mu\ll\pi$. (d) At fixed $\theta$, the latents forget any starting value at the rate of
Theorem~\ref{thm:row}. (e) An exact draw of $\omega$ at a start with fixed $\theta_0$ does not make the start
distribution absolutely continuous, so it adds nothing to (c).
\end{theorem}
\begin{proof}
(a) The latent step is first and $R_i\ge\PG_i/M^*_i$ (Theorem~\ref{thm:row}). (b) Iterating (a),
$S_R^n(\xi,A)\ge\int\epsilon(\theta_0)S_F(\xi,d\xi_1)\epsilon(\theta_1)\cdots S_F(\xi_{n-1},d\xi_n)1_A(\xi_n)$,
positive whenever $S_F^n(\xi,A)>0$. If $S_R$ had a $d$-cycle $D_1,\dots,D_d$ ($d\ge2$, in the sense of
\cite[Sec.~3]{rr}), then for $\xi\in D_i$, $S_F(\xi,D_{i+1}^c)\le S_R(\xi,D_{i+1}^c)/\epsilon(\theta)=0$, a cycle of
$S_F$. (c) \cite[Thm.~4]{rr}: ``If a Markov chain on a state space with countably generated $\sigma$-algebra is
$\varphi$-irreducible and aperiodic, and has a stationary distribution $\pi$, then for $\pi$-a.e. $x$,
$\lim_n\|P^n(x,\cdot)-\pi\|=0$''; the state space (tree shapes times Euclidean spaces) is countably generated;
stationarity is Theorem~\ref{thm:sweep}. For $\mu\ll\pi$, dominated convergence. (d) Theorem~\ref{thm:row}.
(e) The leaves at $\theta_0$ form a point of a continuous space.
\end{proof}

So exactness needs invariance and convergence from the start used; the start itself may be any value (the
installed value, the creation value, or a draw from $q$). What is not proven, here or for today's sampler, is
convergence from every fixed start (Section~\ref{sec:notproven}).

\section{Multinomial}\label{sec:multi}

With categories $k$ and forests $f_k$, category $k$'s rows have shape $n_i$ and $\psi_{ik}=f_{ik}-C_{ik}$,
$C_{ik}=\log\sum_{j\ne k}e^{f_{ij}}$. The build, before forest $k$, draws $\omega_k$ fresh against the current
$f_{-k}$ and then forest $k$ given $\omega_k$: a data-augmentation pair, invariant for $p(f)$.

\begin{theorem}\label{thm:multi}
Keep $\omega_k$ between visits and move it by the step against $\PG(n_i,f_{ik}-C_{ik})$ instead. Under the
extended target $\Pi(f,\omega)=p(f)\prod_k\prod_i\PG(\omega_{ik}\mid n_i,f_{ik}-C_{ik})$ (the one whose latent
conditionals are what the step targets), forest $j$'s draw from $p(f_j\mid f_{-j},\omega_j)$ is not its full
conditional, since $\Pi$'s conditional of $f_j$ carries the factors $\PG(\omega_{ik}\mid n_i,f_{ik}-C_{ik})$,
$k\ne j$, which depend on $f_j$ through $C_{ik}$. The scheme is not exact: in Toy 3 its stationary $f$-marginal is
$0.0061$ from $p$ in $\TV$, while fresh draws are exact.
\end{theorem}
\begin{proof}
The factor computation is direct. Equivalently, every other category's forest step changes $C_{ik}$, hence row
$i$'s conditional, without reading $\omega_k$: the situation of Theorem~\ref{thm:setters}(c) at every row and every
sweep. Toy 3 is in exact arithmetic (Appendix~\ref{app:toys}).
\end{proof}

\section{Theorems and the design elements they cover}\label{sec:map}

\begin{center}\small
\begin{tabular}{>{\raggedright\arraybackslash}p{0.3\textwidth}>{\raggedright\arraybackslash}p{0.62\textwidth}}\toprule
Result & Design element\\\midrule
Prop.~\ref{prop:rev}, Thm.~\ref{thm:row} & the accept-or-keep step itself; its ratio in untilted form; uniform
ergodicity and its rate\\
Lemma~\ref{lem:D}, Props.~\ref{prop:bounded}, \ref{prop:K0}, \ref{prop:Krule}, App.~\ref{app:explicit} & the $K$
rule, the remainder matched at $\psi$, the bounded weight, the rounding of the first rate, the rejection-sampler
research note\\
Lemmas~\ref{lem:contour}--\ref{lem:maj}, Thm.~\ref{thm:dens}, App.~\ref{app:round} & S1's certified densities $f$
and $q$ at every reachable $(b,\psi)$; the pins\\
Prop.~\ref{prop:cf}, Table~\ref{tab:floor} & the tilted, cancellation-free form of the densities and its floor\\
Eqs.~\eqref{eq:tiltupd}, \eqref{eq:lead}, App.~\ref{app:round} & the cache of $\log f$ across sweeps; the lead in
double-double\\
Prop.~\ref{prop:left} & the left series region\\
Lemma~\ref{lem:agg}, Cor.~\ref{cor:tail} & a cheaper certificate for S1 (not built); the cost figures and S1's
timing gate\\
Prop.~\ref{prop:tline} & the vertical line through the saddle point (the earlier large-shape construction)\\
Thm.~\ref{thm:dec}, Props.~\ref{prop:ladder}, \ref{prop:floor}, \ref{prop:fixedtop}, Thm.~\ref{thm:mitro} & S2's
certified decision, refinement, the rounding floor and its tally\\
Section~\ref{sec:ideal}, Lemma~\ref{lem:flip} & what ``exact'' means in floating point, for the step and for
today's sum\\
Prop.~\ref{prop:fresh} & declining a fresh draw on an undecided decision\\
Lemma~\ref{lem:condind}, Thm.~\ref{thm:sweep}, Prop.~\ref{prop:steps} & the step at the sweep's start (latents
first); rows independent given the trees\\
Thm.~\ref{thm:nb} & estimated-$r$ nbinom: fresh where $r$ moved, the step where it stayed\\
Thm.~\ref{thm:rule} & dec-B590's per-row rule\\
Thm.~\ref{thm:setters}, Lemma~\ref{lem:complete}, Cor.~\ref{cor:frac} & marked rows, setters (setData and
setState included), every setter taken at fractional shapes, a fixed set of rows, outer steps of neither kind\\
Thm.~\ref{thm:rule}, the cancel paragraph & the cancel polled between rows\\
Thm.~\ref{thm:starts} & starting points served by the step from the installed value\\
Thm.~\ref{thm:multi} & multinomial keeps the sum\\\bottomrule
\end{tabular}
\end{center}

